\begin{table}[H]
    \centering
    \makebox[\textwidth][c]{%
    \begin{tabular}{|c||c|c|c|}
        \hline
        Anforderung & JUnit XML & TAP & OTR \\
        \hline
        \hline
        \makecell{ermöglicht Ermittlung aller\\ vorhandenen Testcases}
            & \tabNo
            & \tabNo
            & \tabNo \\
        \hline
        eindeutige Identifizierung von Testcases
            & \tabPartial
            & \tabPartial
            & \tabPartial \\
        \hline
        \makecell{sagt für jeden Testcase einzeln aus,\\ ob er ausgeführt wurde oder nicht}
            & \tabPartial
            & \tabPartial
            & \tabPartial \\
        \hline
        strukturiertes Format
            & \tabPartial
            & \tabYes
            & \tabYes \\
        \hline
        weite Verbreitung
            & \tabYes
            & \tabPartial
            & \tabPartial \\
        \hline
    \end{tabular}%
    }
    \vspace{0.5em}
    \begin{tabular}{ll}
        \tabYes & Erfüllt \\
        \tabPartial & In Teilen erfüllt \\
        \tabNo & Nicht erfüllt \\
    \end{tabular}
    \caption{Evaluation der Testausgabeformate nach den in \chapref{Anforderungen Testausgabeformat} definierten Anforderungen}
    \label{tab:evaluation_testausgabeformate}
\end{table}